\documentclass[11pt]{article}
\usepackage{fontspec}
\setmainfont{Times New Roman}
\setsansfont{Arial}
\setmonofont{Consolas}
\usepackage{amsmath,amssymb}
\usepackage{geometry}
\usepackage{microtype}
\geometry{a4paper,margin=3cm}
\setlength{\parindent}{1.25em}
\setlength{\parskip}{0pt}
\emergencystretch=30em
\allowdisplaybreaks
\raggedbottom
\providecommand{\mM}{\mathfrak M}
\providecommand{\mN}{\mathfrak N}
\providecommand{\mA}{\mathfrak A}
\providecommand{\mB}{\mathfrak B}
\providecommand{\mP}{\mathfrak P}
\providecommand{\mD}{\mathfrak D}
\providecommand{\mS}{\mathfrak S}
\providecommand{\ma}{\mathfrak a}
\providecommand{\mb}{\mathfrak b}

\begin{document}

% Унит: Папер 17 Сецтион 10. Дирецт Интерславиц Латин транслатион фром те Герман финал-аудитед соурце слице.

\setcounter{footnote}{19}
\subsection*{\S~10. Обставанье бесконечно многих разкладов групы.}

Ако сумно представјенье пополно редуцибилној групы не јест једнозначно, тогда по Теорему V в всаком представјеньу обставајут најмене две меджу собоју изоморфне подгрупы. Ныне разматрајемо подгрупу, составјену из такој системы изоморфных подгруп,
\[
  \mathfrak H=\mathfrak A_1+\cdots+\mathfrak A_\alpha,
\]
при чем теды \(\mathfrak A_1,\mathfrak A_2,\ldots,\mathfrak A_\alpha\) сут все изоморфне, и тврдимо, же \(\mathfrak H\) потом допушча бесконечно много различных такых разкладов, ако само в области рационалности \(P\) јест бесконечно много ``констант'', то јест величин \(a\), за кторе \(\xi_i a=a\xi_i\). Тут теорем важи чак и тогда, когда о групах \(\mathfrak A_i\) ничто далье не предполагаје се, ни иредуцибилност. Теды важи

\textbf{Теорем VIII.} \textit{Ако \(\mathfrak G\) јест представима како сума конечно многих меджу собоју изоморфных подгруп, тогда тако представјенье јест можно на бесконечно много различных способов, ако само област рационалности \(P\) всебује бесконечно много констант. При том под константами треба разумети таке величины \(a\) области, за кторе \(\xi_i a=a\xi_i\).}

\textbf{Доказ.} Неха теды
\[
  \mathfrak G=\mathfrak A_1+\cdots+\mathfrak A_\alpha,
\]
и \(\mathfrak A_i\) јест изоморфна к \(\mathfrak A_\varkappa\) за \(i,\varkappa=1,\ldots,\alpha\) (\(\alpha\ge2\)). Да бысмо задали определьено изоморфно приредјенье, најпред изделимо једну подгрупу, на примјер \(\mathfrak A_1\). Неха изоморфија меджу \(\mathfrak A_1\) и \(\mathfrak A_i\) буде фиксована через
\begin{equation}
 a_1\sim t_{1i}a_i,\qquad a_i\sim t_{i1}a_1,\qquad t_{11}a_1=a_1
 \quad (i,\varkappa=1,\ldots,\alpha),
\tag{31}
\end{equation}
при чем теды в групе \(\mathfrak A_1\) всак клас остатков јест приредјены сам собе. Понеже \(Ma_1=0\), \(Ma_i=0\), тогда следује \(Mt_{1i}a_i=0\), \(Mt_{i1}a_1=0\); класы \(t_{1i}a_i\), \(t_{i1}a_1\) зато, аналогно јединицам, допушчајут множенье с класами, не само с полиномами. Из дефиниције изоморфије потом следује, же за два приредјене класы \(\mathfrak y\) и \(\bar{\mathfrak y}\), кторе допушчајут тако множенье класов, класы \(\mathfrak r\mathfrak y\) и \(\mathfrak r\bar{\mathfrak y}\) сут знову приредјене. Зато релације \(a_ra_s=0\) \((r\ne s)\) (\(r=1,\ldots,\alpha\); \(s=1,
\ldots,\alpha\)) за \(s=1\) дајут
\begin{equation}
\left\{
\begin{aligned}
 a_r t_{1\varkappa}a_\varkappa&=0 &&(\varkappa=1,
\ldots,\alpha,
\ r\ne1),\\
 \text{а за }s>1\qquad a_r t_{s1}a_1&=0 &&(r\ne s),\\
 \text{теды такоже}\qquad t_{1i}a_i&=a_1t_{1i}a_i,
 & t_{i1}a_1&=a_it_{i1}a_1.
\end{aligned}\right.
\tag{32}
\end{equation}
Далеј из (31), како обобченье (27) и (28), из једнозначности приредјеньа следује
\begin{equation}
 t_{1i}t_{i1}a_1=a_1,
 \qquad
 t_{\varkappa1}t_{1\varkappa}a_\varkappa=a_\varkappa.
\tag{33}
\end{equation}
Као изоморфију меджу \(\mathfrak A_i\) и \(\mathfrak A_\varkappa\) ныне фиксујемо однос, кторы настаје композицијом (31); тако фиксованье јест потребно, понеже једнотливе групы могут допушчати изоморфије в себе самых. Из (31), (32) и (33) тако доставајемо
\begin{equation}
 a_i\sim t_{i1}t_{1\varkappa}a_\varkappa=t_{i\varkappa}a_\varkappa,
 \qquad
 t_{i\varkappa}a_i=a_i,
 \qquad
 a_rt_{s\varkappa}a_\varkappa=0\quad (r\ne s),
\tag{34}
\end{equation}
при чем положено јест \(t_{i1}t_{1\varkappa}=t_{i\varkappa}\). Из (33) и (34) далеј выходи
\[
 t_{i\varkappa}t_{\varkappa l}a_l=t_{i1}t_{1\varkappa}t_{\varkappa1}t_{1l}a_l
 =t_{i1}t_{1l}a_l=t_{il}a_l;
\]
тако јест выдельенье групы \(\mathfrak A_1\) знову одстранено. Изоморфны однос меджу \(\mathfrak A_i\) и \(\mathfrak A_\varkappa\) оставаје тој самы, коју год групу \(\mathfrak A_\varkappa\) употребимо вместо \(\mathfrak A_1\) за дефиницију. Теды \(t_{i\varkappa}\) сут класы, кторе скупај задовольавајут релације
\begin{equation}
 Mt_{i\varkappa}a_\varkappa=0;
 \quad a_rt_{i\varkappa}a_\varkappa=0\ (r\ne i);
 \quad t_{i\varkappa}t_{\varkappa l}a_l=t_{il}a_l;
 \quad t_{\varkappa i}a_i=t_{\varkappa i}t_{i\varkappa}a_\varkappa=a_\varkappa.
\tag{35}
\end{equation}
Те релације говорет точно, же модулы \(\mathfrak M_i\) и \(\mathfrak M_\varkappa\), належече к \(\mathfrak A_i\) и \(\mathfrak A_\varkappa\), сут једнакого вида (сравни дефиницију \S~9); вместо тамошньих \(P\) и \(Q\) тут стојат класы \(t_{i\varkappa}\) односно \(t_{\varkappa i}\), а прве две релације (35), взяте заједно, одповедајут формуле (25) односно (26). Обратно из тих релациј по Теорему VI следује изоморфија.\footnote{I то без примененьа теорема конечности, понеже формулы уже всакократно представајут пресликанье \(\varphi_{i\varkappa}\) и јего обрат \(\varphi_{\varkappa i}\). Понеже \(ra_i=ra_it_{i\varkappa}a_\varkappa t_{\varkappa i}a_i\), из \(ra_i\ne0\) следује такоже \(ra_it_{i\varkappa}a_\varkappa\ne0\).}

Из класов \(t_{i\varkappa}a_\varkappa\) можно ныне составити класы \(b_\rho\), кторе се знову покажут како јединице груп \(\mathfrak B_\rho\), из чего следоват буде представјенье
\[
 \mathfrak G=\mathfrak B_1+\cdots+\mathfrak B_\alpha.
\]
Неха \(\lambda_{i\varkappa}\) сут константы, чија детерминанта \(|\lambda_{i\varkappa}|\ne0\), а \(\lambda^{\varkappa r}\) одповедајучи поддетерминанты, подельене через \(|\lambda_{i\varkappa}|\), тако же обставајут релације
\begin{equation}
\left\{
\begin{aligned}
 \sum_{r=1}^{\alpha}\lambda_{ri}\lambda^{r\varkappa}&=\delta_{i\varkappa}=\begin{cases}0,&i\ne\varkappa,\\1,&i=\varkappa,
 \end{cases}\\
 \sum_{r=1}^{\alpha}\lambda_{ir}\lambda^{\varkappa r}&=\delta_{i\varkappa}.
\end{aligned}\right.
\tag{36}
\end{equation}
Тогда дефинујемо
\begin{equation}
 b_\rho=\sum_{i,\varkappa}\lambda_{\rho i}\lambda^{\rho\varkappa}t_{i\varkappa}a_\varkappa
 \qquad(\rho=1,
\ldots,\alpha).
\tag{37}
\end{equation}
Из (35), (36) и (37) потом следује \(Mb_\rho=0\) и
\begin{equation}
 \sum_{\rho=1}^{\alpha}b_\rho
 =\sum_{i,\varkappa,\rho}\lambda_{\rho i}\lambda^{\rho\varkappa}t_{i\varkappa}a_\varkappa
 =\sum_{i,\varkappa}\delta_{i\varkappa}t_{i\varkappa}a_\varkappa
 =\sum_i t_{ii}a_i=\sum_i a_i=e,
\tag{38}
\end{equation}
далеј
\[
\begin{aligned}
 b_\rho b_\sigma
 &=\sum_{\substack{i,\varkappa\\ \mu,\nu}}
  \lambda_{\rho i}\lambda^{\rho\varkappa}\lambda_{\sigma\mu}\lambda^{\sigma\nu}
  t_{i\varkappa}a_\varkappa t_{\mu\nu}a_\nu \\
 &=\sum_{i,\mu,\nu}\lambda_{\rho i}\lambda^{\rho\mu}\lambda_{\sigma\mu}\lambda^{\sigma\nu}t_{i\nu}a_\nu
 =\delta_{\rho\sigma}\sum_{i,\nu}\lambda_{\rho i}\lambda^{\sigma\nu}t_{i\nu}a_\nu
 =\begin{cases}0&(\rho\ne\sigma),\\ b_\rho&(\rho=\sigma),\end{cases}
\end{aligned}
\]
\((\rho,\sigma=1,
\ldots,\alpha)\). Представјенье \(\mathfrak r e=\mathfrak r b_1+\cdots+\mathfrak r b_\alpha\) јест теды адитивны разклад групы \(\mathfrak G\). При том групы \(\mathfrak B_\rho\) сут различне од груп \(\mathfrak A_\sigma\), ако скоро \(\lambda_{i\varkappa}\) не творе само диагоналну матрицу. Бо ако формујемо \(a_\sigma b_\rho\), по (35) доставајемо:
\[
 a_\sigma b_\rho=
 \sum_{i,\varkappa}\lambda_{\rho i}\lambda^{\rho\varkappa}a_\sigma t_{i\varkappa}a_\varkappa
 =\lambda_{\rho\sigma}
 \sum_{\varkappa}\lambda^{\rho\varkappa}t_{\sigma\varkappa}a_\varkappa\ne0
 \quad\text{за }\lambda_{\rho\sigma}\ne0,
\]
понеже иначе в последньој суме мусил бы изчезнути всакы поједины счлен, а то не става се. Ако теды при фиксованом \(\sigma\) обставајут најмене два \(\rho\), за кторе \(\lambda_{\rho\sigma}\ne0\), тогда \(\mathfrak A_\sigma\) не може быти идентична с никојом \(\mathfrak B_\rho\).

Но групы \(\mathfrak B_\rho\) сут такоже изоморфне меджу собоју и к групам \(\mathfrak A_\sigma\). За то најпред показујемо обставанье класов \(r_{\rho\sigma}a_\sigma\) и \(r^{\sigma\rho}b_\rho\), кторе посредничајут изоморфију \(\mathfrak A_\sigma\) и \(\mathfrak B_\rho\) (\(\rho,\sigma=1,
\ldots,\alpha\)). Доказујемо именно, же за класы
\[
 r_{\rho\sigma}a_\sigma=\sum_i\lambda_{\rho i}t_{i\sigma}a_\sigma,
 \qquad
 r^{\sigma\rho}=\sum_\varkappa\lambda^{\rho\varkappa}t_{\sigma\varkappa}a_\varkappa
 =r^{\sigma\rho}b_\rho
\]
релације ``једнакого вида'' меджу модулами
\[
 (\mathfrak M,b_1+\cdots+b_{\rho-1}+b_{\rho+1}+\cdots+b_\alpha)
\]
и
\[
 (\mathfrak M,a_1+\cdots+a_{\sigma-1}+a_{\sigma+1}+\cdots+a_\alpha)
\]
сут изполньене:
\begin{equation}
\left\{
\begin{aligned}
 b_\varkappa r_{\rho\sigma}a_\sigma&=0 &&(\varkappa\ne\rho),&
 r^{\sigma\rho}r_{\rho\sigma}a_\sigma&=a_\sigma,\\
 a_\varkappa r^{\sigma\rho}b_\rho&=0 &&(\varkappa\ne\sigma),&
 r_{\rho\sigma}r^{\sigma\rho}b_\rho&=b_\rho,
\end{aligned}\right.
\tag{39}
\end{equation}
из чего, како при (35), следује изоморфија посредством приредјеньа
\begin{equation}
 b_\rho\sim r_{\rho\sigma}a_\sigma,
 \qquad
 a_\sigma\sim r^{\sigma\rho}b_\rho.
\tag{40}
\end{equation}
В самој дејствителности выходи:
\[
\begin{aligned}
 b_\varkappa r_{\rho\sigma}a_\sigma
 &=\sum_{\mu,\nu,i}\lambda_{\varkappa\mu}\lambda^{\varkappa\nu}t_{\mu\nu}a_\nu\lambda_{\rho i}t_{i\sigma}a_\sigma\\
 &=\sum_{\mu,i}\lambda_{\varkappa\mu}\lambda^{\varkappa i}\lambda_{\rho i}t_{\mu\sigma}a_\sigma
 =\delta_{\varkappa\rho}\sum_\mu\lambda_{\varkappa\mu}t_{\mu\sigma}a_\sigma
 =\delta_{\varkappa\rho}r_{\rho\sigma}a_\sigma=0\quad(\varkappa\ne\rho),
\end{aligned}
\]
далеј
\[
 r^{\sigma\rho}r_{\rho\sigma}a_\sigma
 =\sum_{\nu,i}\lambda^{\rho\nu}t_{\sigma\nu}a_\nu\lambda_{\rho i}t_{i\sigma}a_\sigma
 =\sum_i\lambda^{\rho i}\lambda_{\rho i}a_\sigma=a_\sigma.
\]
Такоже обратно находимо
\[
 a_\varkappa r^{\sigma\rho}b_\rho
 =a_\varkappa\sum_\mu\lambda^{\rho\mu}t_{\sigma\mu}a_\mu=0
 \quad(\varkappa\ne\sigma),
\]
далеј
\[
 r_{\rho\sigma}r^{\sigma\rho}b_\rho
 =\sum_{i,\mu}\lambda_{\rho i}t_{i\sigma}a_\sigma\lambda^{\rho\mu}t_{\sigma\mu}a_\mu
 =\sum_{i,\mu}\lambda_{\rho i}\lambda^{\rho\mu}t_{i\mu}a_\mu=b_\rho.
\]
Из тут доказаних релациј једнакого вида ныне следује (знову без теорема конечности) изоморфија \(\mathfrak A_\sigma\) и \(\mathfrak B_\rho\). Изоморфија меджу \(\mathfrak B_\rho\) и \(\mathfrak B_\tau\) находи се одтуд композицијом; за то јошче дајемо формулы. Положимо
\[
 \ell_{\rho\tau}b_\tau=r_{\rho\sigma}r^{\sigma\tau}b_\tau
 =\sum_{i,j}\lambda_{\rho i}t_{i\sigma}a_\sigma\lambda^{\tau j}t_{\sigma j}a_j
 =\sum_{i,j}\lambda_{\rho i}\lambda^{\tau j}t_{ij}a_j;
\]
тогда приредјенье гласи \(b_\rho\sim\ell_{\rho\tau}b_\tau\). I тут можно релације једнакого вида знову формално доказати; \(\ell_{\rho\tau}\), поставльене на место \(t_{i\varkappa}\), изполньајут релације (35). В самој дејствителности
\[
 b_\sigma\ell_{\rho\tau}b_\tau=b_\sigma r_{\rho\sigma}r^{\sigma\tau}b_\tau=0\quad(\sigma\ne\rho),
\]
\[
 \ell_{\rho\sigma}\ell_{\sigma\tau}b_\tau=r_{\rho\varkappa}r^{\varkappa\sigma}r_{\sigma\lambda}r^{\lambda\tau}b_\tau
 =r_{\rho\varkappa}a_\varkappa r^{\varkappa\tau}b_\tau=\ell_{\rho\tau}b_\tau.
\]

Тым Теорем VIII јест доказан, понеже треба само избирати \(\lambda_{i\varkappa}\) на бесконечно много способов тако, же оны не творе диагоналну матрицу и не могут быти трансформоване једна в другу через диагоналну матрицу.

Последство Теорема VIII, в повезце с Теоремом V односно VII, јест

\textbf{Теорем IX.}\footnote{Сравни Лоевы, с. 106--107.} \textit{Ако модул \(\mathfrak M\) јест пополно редуцибилны и јест најмене обче многократно простых модулов \(\mathfrak P_1,
\ldots,
\mathfrak P_k\), кторе творе тотално взаимно просту систему, и ако област рационалности \(P\) всебује бесконечно много констант, тогда нужна и достаточна условја за то, же \(\mathfrak M\) јест представимы на бесконечно много различных способов како најмене обче многократно простых модулов, јест та, же најмене два из модулов \(\mathfrak P_1,
\ldots,
\mathfrak P_k\) сут једнакого вида.}

Ныне хочемо вывести обче критерије за то, же модул \(\mathfrak M\) јест делителны через бесконечно много простых модулов, и за то најпред доказујемо следујучи

\textbf{Помочны теорем.}\footnote{Сравни Лоевы, с. 96.} \textit{Ако пополно редуцибилны модул \(\mathfrak M=[\mathfrak P_1,
\ldots,
\mathfrak P_k]\) јест делителны через просты модул \(\mathfrak P'\), различны од всех \(\mathfrak P_i\), тогда најмене два из \(\mathfrak P_i\) сут једнакого вида.}

Најпред из \(\mathfrak P_1,
\ldots,
\mathfrak P_k\) избирамо тотално взаимно просту систему, чртајучи по реду всакы из тих простых модулов, кторы входи в најмене обче многократно јошче не изчртаних. Зато ныне неха буде предположено, же \(\mathfrak P_1,
\ldots,
\mathfrak P_k\) саме творе тотално взаимно просту систему. Неха \(a'\) јест клас модула \(\mathfrak M\), приредјены јединичному класу мод \(\mathfrak P'\). Тогда
\[
 Fa'=Fa'a_1+\cdots+Fa'a_k
\]
за всакы полином \(F\). Из продуктов \(a'a_i\) сут најмене два различне од нуле, на примјер \(a'a_1\) и \(a'a_2\); бо ако бы, на примјер, само \(a'a_1\ne0\), тогда бы было \(Fa'=Fa'a_1\), и класы остатков \(Fa'\) творили бы подгрупу \(\mathfrak A_1\), были бы теды, понеже \(\mathfrak A_1\) јест проста група, идентичне с \(\mathfrak A_1\), тако же бы се и модулы \(\mathfrak P'\) и \(\mathfrak P_1\) совпадали. По начину закльучиваньа из \S~8 из того следује изоморфија \(\mathfrak A'\) с \(\mathfrak A_1\) и \(\mathfrak A_2\).

Неха ныне \(\mathfrak M\) јест либовольны модул. Тогда разматрајемо најмене обче многократно \(\mathfrak N\) всех простых модулов, через кторе \(\mathfrak M\) јест делителны.\footnote{Ако јих јест бесконечно много, тогда тому представјеньу по Теорему IV не може одповедати никакы адитивны разклад групы остатков. Сравни впрочем примјер \S~12, 4.} Тада сут можне два случаји. Или \(\mathfrak N\) не јест представимы како најмене обче многократно конечно многих простых модулов; в том случају \(\mathfrak N\), теды такоже \(\mathfrak M\), имаје бесконечно много простых модулов како делителье. Или \(\mathfrak N\) јест представимы како најмене обче многократно конечно многих простых модулов; тогда \(\mathfrak N\) называје се највечши пополно редуцибилны модул од \(\mathfrak M\) (в повезце с одповедајучо творбо у Лоевы). Ако медју тими конечно многими простыми модулами два сут једнакого вида, тогда по Теорему IX модул \(\mathfrak N\), и последовательно \(\mathfrak M\), јест делителны через бесконечно много простых модулов. Ако последнье обратно важи, тогда обставаје најмене једин просты делитель од \(\mathfrak N\), кторы јест различны од всех тих конечно многих, како чије најмене обче многократно \(\mathfrak N\) јест представимы. По помочном теорему тогда \(\mathfrak N\), теды такоже \(\mathfrak M\), имаје два просте делителье једнакого вида. Тым доказали јесмо следујучи теорем:

\textbf{Теорем X.} \textit{Ако област рационалности \(P\) всебује бесконечно много констант, тогда нужна и достаточна условја за то, же модул јест делителны через бесконечно много различных простых модулов, јест та, же или не обставаје највечши пополно редуцибилны модул, или, ако такы обставаје, же он јест делителны через најмене два различне просте модулы једнакого вида.}

\end{document}

